\section{Welfare Analysis} \label{sec:mvpf}

\subsection{The Capitalization Gap}

The combination of a near-zero immediate price response and a large eventual cumulative decline implies that households who voted for the tax cut accepted a wealth loss substantially larger than the present value of the savings they realized. Before turning to a fully structural welfare calculation, it is useful to express this {\it capitalization gap} in transparent dollar terms.

The mean median sale price in our sample is approximately \$166{,}000. The estimated cumulative ten-year decline in median sale price following a narrow rejection of a renewal road levy is \$15{,}350, or roughly nine percent. The average per-household annual tax savings from rejecting the renewal is approximately \$76 (Table \ref{tab:levy_stats}), and the typical levy duration before re-vote is five years. Capitalizing the savings as a perpetuity at the user-cost-of-housing discount rate $r \approx 0.04$ yields a present value of approximately \$1{,}900 per household.\footnote{We use the perpetuity approximation rather than the five-year levy horizon because, in equilibrium, the typical jurisdiction re-votes on a road levy every five years and the rejection rate at any future vote is plausibly similar to the contemporaneous rate. A five-year horizon would yield a lower-bound estimate of approximately \$340; the qualitative magnitude of the capitalization gap is invariant to this choice.}

The implied {\it capitalization gap} is therefore $\$15{,}350 / \$1{,}900 \approx 8$: every dollar of present-value tax savings accepted by the median household is associated with eight dollars of capitalized housing-wealth loss. This is the headline ratio we report in the introduction. The gap is large precisely because the tax cut is small relative to the present value of the amenity stream it finances, and because the maintenance shortfall compounds nonlinearly through the road-capital law of motion in Equation \ref{eq:roads}.

\subsection{Marginal Value of Public Funds}

To translate the capitalization gap into a standardized welfare metric we follow \cite{hendren2020unified} and \cite{finkelstein2020welfare} and compute the marginal value of public funds (MVPF), defined as the ratio of household willingness to pay (WTP) for the tax cut to the government's net cost of providing it:
\begin{equation}
\text{MVPF} \;=\; \frac{\text{WTP}}{\text{Net Cost to Government}}
\;=\; \frac{\text{Tax Savings} - \text{House Value Decline}}{\text{Direct Long Run Cost} + \text{Fiscal Externality}}.
\end{equation}
The components are:
\begin{itemize}
    \item \textbf{Direct Long Run Cost} is the per-household revenue loss to the government over the lifetime of the levy reduction.
    \item \textbf{Fiscal Externality} is the additional revenue loss arising from the decline in the property-tax base induced by capitalized house-price declines.
    \item \textbf{Tax Savings} is the discounted value of the household's foregone road-maintenance tax payments.
    \item \textbf{House Value Decline} is the average yearly capitalized loss in house values over the period of analysis.
\end{itemize}

Plugging in our empirical estimates yields an MVPF of $-1.4$: for every dollar of revenue forgone by the government, households incur a net welfare loss of \$1.4 after accounting for both direct tax savings and the fiscal externality. This is a welfare-destroying policy in the sense of \citet{hendren2020unified}. While they do not report a comparable infrastructure-specific MVPF, their estimates for adult-targeted policies fall in the range $[0.5,\, 2]$, placing the road-tax-cut MVPF well outside the support of welfare-improving fiscal policies they examine.

The MVPF is, however, a {\it conservative} measure of the welfare cost of the policy. It omits two channels that the structural model in Section \ref{sec:mech} makes explicit. First, the labor-supply distortion induced by elevated $\phi(G_t)$ (commuting disutility) reduces output by approximately three percent in steady state, generating a direct deadweight loss not captured by housing capitalization alone. Second, the consumption-smoothing benefit of higher disposable income in the immediate post-vote period is partly offset by the long-run drop in steady-state consumption that the model traces. A full structural welfare calculation accounting for these two channels yields a more negative MVPF; we report $-1.4$ as the headline number because it is invariant to the parameters of the production function and to the elasticity of housing supply.

\subsection{Behavioral Welfare Adjustment}

The MVPF computation above presumes that household choices reveal preferences --- that voters who supported the levy cut did so with a correct understanding of its consequences. As discussed in Section \ref{sec:fiscal_illusion}, three candidate mechanisms can rationalize the observed delayed-capitalization pattern, and they differ in whether this revealed-preference presumption is appropriate.

\paragraph{The buyer-side benchmark (Mechanisms A and B).}
Under the belief-updating microfoundation of Section \ref{sec:belief_updating}, the revealed-preference presumption is violated by construction. At the moment of the vote, the median household's posterior $\hat{G}_t$ is anchored at its pre-election baseline, and the perceived utility cost of the tax cut is therefore close to zero. The realized cost --- the eventual capitalization of the maintenance shortfall --- accrues only after $\hat{G}_t$ has adjusted, by which point the vote is irreversible. Following \cite{chetty2009salience}, we compute a counterfactual MVPF using the {\it experienced-utility} cost of the tax cut rather than the cost implied by household choices. We evaluate the policy at the perfect-foresight solution of the model, in which households fully internalize the future amenity decline at $t = 0$. The PF counterfactual implies that no rational household with $\hat{G}_t = G_t$ at the vote date would have supported the tax cut. The wedge between the choice-based MVPF and this experience-based MVPF is the structural counterpart to the capitalization gap of Section \ref{sec:fiscal_illusion}. Quantitatively, the experience-based MVPF lies in the range $[-2.1,\, -1.6]$ across the calibrated range of $\bar{\sigma}$.

\paragraph{The voter-side benchmark (Mechanism C: voter myopia).}
Under voter myopia, the revealed-preference presumption must be interpreted carefully. The choice-based MVPF of $-1.4$ accurately summarizes what voters chose at the moment of the vote; whether that choice reflects their {\it true} long-horizon preferences depends on the nature of the myopia. If the myopia is cognitive --- voters did not fully comprehend the magnitude or timing of the future amenity decline --- then an informed counterfactual would yield a more negative MVPF closer to the experience-based estimate, because voters would have voted against the cut under full information. If the myopia is preference-based --- voters fully understood the future loss but rationally discounted it because of a steep subjective discount factor or a short expected housing tenure --- then the choice-based MVPF approximately captures voter welfare, because the voter would have made the same decision even with full information.

\paragraph{Range of MVPF estimates.}
Because voters likely fall on a spectrum between purely-cognitive and purely-preference-based myopia, the welfare-relevant MVPF lies somewhere between the experience-based bound (as low as $-2.1$) and the choice-based bound ($-1.4$). The qualitative conclusion --- that the road-maintenance tax cut is welfare-reducing --- is robust across the range. The quantitative conclusion --- and the implied benefit of policy intervention --- depends on (i) the relative weight of buyer-side versus voter-side frictions and (ii) within the voter-side friction, the relative weight of cognitive versus preference-based myopia.

\paragraph{Implications for the optimal policy lever.}
The distinction has direct policy bite. If buyer-side frictions dominate, information-disclosure policies (e.g., mandatory pre-vote road-condition assessments, in the spirit of restaurant-grading and energy-efficiency disclosure) deliver large welfare gains by reducing $\bar{\sigma}$. If voter-side cognitive myopia dominates, the same disclosure lever applies but operates through a different channel --- it educates voters about the eventual cost rather than buyers about the current state. If voter-side preference-based myopia dominates (a steep subjective discount or a short tenure horizon), disclosure delivers smaller gains; the welfare-improving lever is instead institutional --- pre-commitment devices that bind future maintenance funding (e.g., constitutionally-protected dedicated revenue streams, automatic indexation of road levies to maintenance cost inflation) so that the temporally-distant cost is not bypassed by a myopic voter at any single ballot.

\subsection{Policy Implications}

Two policy levers are suggested by the welfare analysis above. First, mandatory pre-vote disclosure of road-condition assessments would reduce the noise variance $\sigma_t^2$ in Equation \ref{eq:signal} at precisely the moment voters most need it. Mechanisms with this intent --- restaurant-grading systems in New York City, energy-efficiency disclosure at sale, water-quality reporting --- have been shown to alter consumer behavior at modest implementation cost. The fire-protection findings of \citet{brasington2024fire} provide a particularly close analog: a limited-awareness mechanism for capitalization of public services that attenuates when awareness rises. Second, unbundling road-maintenance levies from other items on the ballot would reduce the cognitive cost emphasized in the limited-attention microfoundation of Section \ref{sec:fiscal_illusion}.

The model-based welfare calculation suggests that even modest reductions in $\bar{\sigma}$ would deliver substantial welfare gains. A counterfactual in which $\bar{\sigma}$ is halved --- corresponding, roughly, to a doubling of pre-vote attention --- reduces the steady-state capitalization gap by approximately forty percent. The full counterfactual analysis is reported in Appendix \ref{sec:appxdmodel}.
